\documentclass[10pt,a4paper]{amsart}
\usepackage[margin=19mm]{geometry}
\usepackage{amsmath,amssymb,mathtools}
\usepackage[scr=rsfs,cal=euler]{mathalfa}
\usepackage{xfrac}
\usepackage[unicode,hidelinks,bookmarks=false]{hyperref}
\newcommand{\ie}{i.e.\ }
\newcommand{\cf}{cf.\ }
\newcommand{\resp}{respectively\ }
\newcommand{\Subsection}{\subsection}
\providecommand{\nobreakdash}{\nobreak\hbox{-}}
\setlength{\emergencystretch}{2em}
\multlinegap0pt
\usepackage{algorithm}\usepackage{algpseudocode}
\usepackage{amssymb}
\usepackage{float}
\newtheorem{lemma}{Lemma}
\newcommand{\PP}{\mathbb{P}}
\newcommand{\cE}{\mathcal{E}}
\newcommand{\cbr}[1]{\left\{#1\right\}}
\newcommand{\floor}[1]{\left\lfloor#1\right\rfloor}
\newcommand{\rbr}[1]{\left(#1\right)}
\title{\normalfont Simultaneous Pointwise Majorization for Mixed Tail Processes with Applications in Gaussian Chaos and Ergodic Diffusions}
\author{Haichen Hu, David Simchi-Levi}
\date{Source: arXiv:2609.01576v2 (CC BY 4.0)}
\begin{document}
\maketitle

\noindent\emph{Source and licence.} \normalfont Simultaneous Pointwise Majorization for Mixed Tail Processes with Applications in Gaussian Chaos and Ergodic Diffusions. Version arXiv:2609.01576v2 (CC BY 4.0), distributed by arXiv under the Creative Commons Attribution 4.0 International licence, http://creativecommons.org/licenses/by/4.0/, as recorded in the arXiv licence metadata for this exact version and shown on the abstract page https://arxiv.org/abs/2609.01576v2. Full licence notice: \texttt{licenses/arxiv-2609.01576-CC-BY-4.0.txt}.\par\smallskip

\noindent\textit{Editorial scope.} Counted unit: the article's lemma lem:mixed\_tail\_finite\_pointwise (source0.tex lines 3501--3527). The article's complete original proof of it is delivered with it (source0.tex lines 3529--3901, one \texttt{proof} environment of 373 lines). No mathematics is written, repaired or re-derived: the statement and the proof are the article's own lines, in the article's own notation, and no retained line is substituted or re-flowed.  No further source-local context is needed: every cross-reference of every retained line resolves inside the two delivered spans.  Nothing was omitted from any retained span, so the package carries no omission record.  No retained line contains a citation, so the wrapper carries no bibliography and no bibliography entry.  No retained line contains a figure environment, an image-inclusion command or drawing code, so no image is delivered and the package declares no asset dependency.  Editorial formatting only, in addition: an AMS \texttt{amsart} preamble that restates the article's own declarations and macros used by the retained lines, namely the environments \texttt{lemma} on separate counters, together with the standard packages those lines need.  The retained environments print in this document as Lemma 1 in order of appearance, whereas the article prints the same items with the numbering of its own full document.  The complete native source of the article as distributed in the arXiv e-print for this exact version is bundled unchanged as \texttt{sources/209126/source0.tex}.
\begin{lemma}
\label{lem:mixed_tail_finite_pointwise}
Let $K\subseteq T$ be finite and contain $t_0$. For every
$j\in\cbr{1,2}$, let $\mathcal T_j=(T_{j,k})_{k\ge0}$ be an admissible sequence of subsets of $K$ such that $T_{j,0}=\cbr{t_0}$ and $T_{j,k}=K$ for all sufficiently large $k$. We define
\begin{align*}
G_j(t)
:=
\sum_{k\ge0}
2^{k/\alpha_j}d_j(t,T_{j,k}),
\ t\in K.
\end{align*}
There exists a finite constant $A_{\alpha_1,\alpha_2}$, depending only
on $\alpha_1$ and $\alpha_2$, such that, for every $u\ge1$, with
probability at least $1-e^{-u}$, simultaneously for every $t\in K$,
\begin{align*}
\|Z_t\|
\le
A_{\alpha_1,\alpha_2}
\cbr{
G_1(t)+G_2(t)
+
u^{1/\alpha_1}v_1(t)
+
u^{1/\alpha_2}v_2(t)
}.
\end{align*}
\end{lemma}

\begin{proof}[Proof of Lemma
\ref{lem:mixed_tail_finite_pointwise}]
For $j\in\cbr{1,2}$, $k\ge0$, and $t\in K$, define the nearest point map $p_{j,k}(t)\in T_{j,k}$ such that
\begin{align*}
d_j(t,p_{j,k}(t))
=
d_j(t,T_{j,k}).
\end{align*}
Such a point exists because $T_{j,k}$ is a nonempty finite set. When
$T_{j,k}=K$, we choose $p_{j,k}(t):=t$.

For every $n\ge2$, define an equivalence relation on $K$ by
\begin{align*}
s\sim_n t\ \text{iff}\ p_{j,k}(s)=p_{j,k}(t)
\end{align*}
for every $j\in\cbr{1,2}$ and every $0\le k\le n-2$. Let
$\mathcal A_n$ be the corresponding partition of $K$, and we write
$\mathcal A_n(t)$ for the cell containing $t$. Equivalently, we define
\begin{align*}
P_n(t)
:=
\rbr{
p_{1,0}(t),\ldots,p_{1,n-2}(t),
p_{2,0}(t),\ldots,p_{2,n-2}(t)
}.
\end{align*}
Then two points lie in the same cell of $\mathcal A_n$ precisely when
they have the same image under $P_n$.

The partitions are nested. Indeed, if $s\sim_n t$, then the nearest
point labels of $s$ and $t$ agree at all levels $0\le k\le n-2$, and
therefore they also agree at all levels $0\le k\le n-3$. Hence,
\begin{align*}
s\sim_n t
\Rightarrow
s\sim_{n-1}t.
\end{align*}
Thus, $\mathcal A_n$ refines $\mathcal A_{n-1}$ for every $n\ge3$.

We next bound the number of cells. By admissibility of the two
sequences, we have that
\begin{align*}
|\mathcal A_n|\le
\prod_{j=1}^2
\prod_{k=0}^{n-2}
|T_{j,k}|\le
\prod_{j=1}^2
\prod_{k=0}^{n-2}
2^{2^k}=
2^{2\sum_{k=0}^{n-2}2^k}=
2^{2(2^{n-1}-1)}=
2^{2^n-2}\le
2^{2^n}.
\end{align*}

For every $n\ge2$, define the pointwise combined approximation cost
\begin{align*}
R_n(t)
:=
2^{n/\alpha_1}d_1(t,T_{1,n-2})
+
2^{n/\alpha_2}d_2(t,T_{2,n-2}).
\end{align*}
For every cell $A\in\mathcal A_n$, choose one point $r_n(A)$ such that
\begin{align*}
r_n(A)\in\operatorname*{argmin}_{s\in A}R_n(s),
\end{align*}
and define $\pi_n(t):=r_n(\mathcal A_n(t))$.
Thus, $\pi_n$ is constant on every cell of $\mathcal A_n$.

Since $t$ and $\pi_n(t)$ belong to the same cell, they have the same
nearest-point label at level $n-2$. Therefore, we have that
\begin{align*}
p_{j,n-2}(t)
=
p_{j,n-2}(\pi_n(t)),
\qquad j\in\cbr{1,2}.
\end{align*}
By the triangle inequality, we obtain
\begin{align*}
d_j(t,\pi_n(t))\le
d_j(t,p_{j,n-2}(t))
+
d_j(p_{j,n-2}(\pi_n(t)),\pi_n(t))=
d_j(t,T_{j,n-2})
+
d_j(\pi_n(t),T_{j,n-2}).
\end{align*}
Multiplying by $2^{n/\alpha_j}$ and summing over
$j\in\cbr{1,2}$, we obtain
\begin{align*}
&2^{n/\alpha_1}d_1(t,\pi_n(t))
+
2^{n/\alpha_2}d_2(t,\pi_n(t))\le
R_n(t)+R_n(\pi_n(t)).
\end{align*}
Since $\pi_n(t)$ is the minimal distance, we have
\begin{align*}
R_n(\pi_n(t))
\le
R_n(t).
\end{align*}
Consequently,
\begin{align}
2^{n/\alpha_1}d_1(t,\pi_n(t))
+
2^{n/\alpha_2}d_2(t,\pi_n(t))
\le
2R_n(t).
\label{eq:mixed_common_representative_bound}
\end{align}

Define $a_{\alpha_1,\alpha_2}
:=
\max\cbr{2^{1/\alpha_1},2^{1/\alpha_2}}$.
For every $n\ge3$, the triangle inequality gives
\begin{align*}
d_j(\pi_n(t),\pi_{n-1}(t))
\le
d_j(\pi_n(t),t)
+
d_j(t,\pi_{n-1}(t)).
\end{align*}
Using \eqref{eq:mixed_common_representative_bound} at level $n$, we
obtain $\sum_{j=1}^2
2^{n/\alpha_j}d_j(\pi_n(t),t)
\le
2R_n(t)$.

At level $n-1$, we use $2^{n/\alpha_j}
\le
a_{\alpha_1,\alpha_2}
2^{(n-1)/\alpha_j}$
and obtain that
\begin{align*}
&\sum_{j=1}^2
2^{n/\alpha_j}d_j(t,\pi_{n-1}(t))\le
a_{\alpha_1,\alpha_2}
\sum_{j=1}^2
2^{(n-1)/\alpha_j}d_j(t,\pi_{n-1}(t))\le
2a_{\alpha_1,\alpha_2}R_{n-1}(t).
\end{align*}
Therefore, combining the inequalities above, we have
\begin{align}
&2^{n/\alpha_1}
d_1(\pi_n(t),\pi_{n-1}(t))
+
2^{n/\alpha_2}
d_2(\pi_n(t),\pi_{n-1}(t))\le
2R_n(t)
+
2a_{\alpha_1,\alpha_2}R_{n-1}(t).
\label{eq:mixed_common_parent_edge_bound}
\end{align}

Because the nearest-point maps are eventually the identity, the
partitions $\mathcal A_n$ are eventually the singleton partition of
$K$. Hence, for every $t\in K$, we have $\pi_n(t)=t$ for all sufficiently large $n$.

For every $n\ge3$, define the set of distinct parent-child edges by
\begin{align*}
E_n
:=
\cbr{
(\pi_n(t),\pi_{n-1}(t)):
t\in K
}.
\end{align*}
Since $\mathcal A_n$ refines $\mathcal A_{n-1}$, every cell of
$\mathcal A_n$ is contained in a unique cell of
$\mathcal A_{n-1}$. Both $\pi_n$ and $\pi_{n-1}$ are constant on a
cell of $\mathcal A_n$. Thus, every cell of $\mathcal A_n$ determines
at most one edge in $E_n$, and
\begin{align}
|E_n|
\le
|\mathcal A_n|
\le
2^{2^n}.
\label{eq:mixed_common_edge_cardinality}
\end{align}

We first suppose that $u\ge4$ and define $\ell:=\floor{\log_2u}$. Then, we have
\begin{align*}
\ell\ge2,
\ 
2^\ell\le u<2^{\ell+1}.
\end{align*}
We now apply the mixed-tail increment bound with parameter $16u$ to all
distinct coarse edges $\cbr{(\pi_\ell(t),t_0),
\ t\in K}$. The number of distinct points $\pi_\ell(t)$ is at most
$|\mathcal A_\ell|$. Hence, by a union bound, we have
\begin{align*}
&\PP\left\{
\exists t\in K\ \text{s.t.} \|Z_{\pi_\ell(t)}-Z_{t_0}\|
>
(16u)^{1/\alpha_1}
d_1(\pi_\ell(t),t_0)
+
(16u)^{1/\alpha_2}
d_2(\pi_\ell(t),t_0)
\right\}\\
\le&2|\mathcal A_\ell|e^{-16u}\le
2\exp\cbr{(\log2)2^\ell-16u}\le
2e^{-15u}.
\end{align*}
For every $n>\ell$, apply the mixed-tail increment bound with parameter
$16\cdot2^n$ to every edge in $E_n$. By
\eqref{eq:mixed_common_edge_cardinality} and another union bound,
\begin{align*}
&\PP\left\{\exists (s,r)\in E_n\ \text{s.t.}\|Z_s-Z_r\|
>
(16\cdot2^n)^{1/\alpha_1}d_1(s,r)
+
(16\cdot2^n)^{1/\alpha_2}d_2(s,r)
\right\}\\
&\le
2|E_n|e^{-16\cdot2^n}\le
2\exp\cbr{-(16-\log2)2^n}\le
2e^{-15\cdot2^n}.
\end{align*}
Since $2^{\ell+1}>u$, by direct algebra, we have
\begin{align*}
\sum_{n>\ell}2e^{-15\cdot2^n}\le
2\sum_{r\ge1}e^{-15u2^{r-1}}\le
2\sum_{r\ge1}e^{-15ur}=
\frac{2e^{-15u}}{1-e^{-15u}}.
\end{align*}
For $u\ge4$,  we have that $2e^{-15u}
\le
\frac12e^{-u},
\ 
\frac{2e^{-15u}}{1-e^{-15u}}
\le
\frac12e^{-u}$. 

Therefore, outside an event $\cE_u$ of probability at most $e^{-u}$,
all the preceding coarse and fine increment bounds hold
simultaneously. We now work conditioning on this event $\cE_u$. 

Since $\pi_n(t)=t$ for all sufficiently
large $n$ and $Z_{t_0}=0$, we may telescope and obtain, simultaneously
for every $t\in K$,
\begin{align*}
\|Z_t\|
&\le
\|Z_{\pi_\ell(t)}-Z_{t_0}\|
+
\sum_{n>\ell}
\|Z_{\pi_n(t)}-Z_{\pi_{n-1}(t)}\|\\
&\le
\sum_{j=1}^2
(16u)^{1/\alpha_j}
d_j(\pi_\ell(t),t_0)+
\sum_{n>\ell}
\sum_{j=1}^2
(16\cdot2^n)^{1/\alpha_j}
d_j(\pi_n(t),\pi_{n-1}(t)).
\end{align*}

We bound these two terms separately. For the first term, by the
triangle inequality, we have
\begin{align*}
d_j(\pi_\ell(t),t_0)\le
d_j(\pi_\ell(t),t)+d_j(t,t_0)=
d_j(t,\pi_\ell(t))+v_j(t).
\end{align*}
Define $c_{\alpha_1,\alpha_2}
:=
\max\cbr{16^{1/\alpha_1},16^{1/\alpha_2}}$. Since $u<2^{\ell+1}$, we have
\begin{align*}
u^{1/\alpha_j}
<
2^{1/\alpha_j}2^{\ell/\alpha_j}
\le
a_{\alpha_1,\alpha_2}2^{\ell/\alpha_j}.
\end{align*}
It follows from
\eqref{eq:mixed_common_representative_bound} that
\begin{align*}
\sum_{j=1}^2
(16u)^{1/\alpha_j}
d_j(t,\pi_\ell(t))\le
c_{\alpha_1,\alpha_2}
a_{\alpha_1,\alpha_2}
\sum_{j=1}^2
2^{\ell/\alpha_j}d_j(t,\pi_\ell(t))\le
2c_{\alpha_1,\alpha_2}^{(16)}
a_{\alpha_1,\alpha_2}R_\ell(t).
\end{align*}
The anchor part satisfies that
\begin{align*}\sum_{j=1}^2
(16u)^{1/\alpha_j}v_j(t)\le
c_{\alpha_1,\alpha_2}
\cbr{
u^{1/\alpha_1}v_1(t)
+
u^{1/\alpha_2}v_2(t)
}.
\end{align*}

For the term $\sum_{n>\ell}
\sum_{j=1}^2
(16\cdot2^n)^{1/\alpha_j}
d_j(\pi_n(t),\pi_{n-1}(t))$, by inequality
\eqref{eq:mixed_common_parent_edge_bound}, we have that
\begin{align*}
\sum_{n>\ell}
\sum_{j=1}^2
(16\cdot2^n)^{1/\alpha_j}
d_j(\pi_n(t),\pi_{n-1}(t))\le&
c_{\alpha_1,\alpha_2}
\sum_{n>\ell}
\cbr{
2R_n(t)
+
2a_{\alpha_1,\alpha_2}R_{n-1}(t)
}\\
\le&
2c_{\alpha_1,\alpha_2}
\rbr{1+a_{\alpha_1,\alpha_2}}
\sum_{n\ge2}R_n(t).
\end{align*}
Combining these inequalities together, we obtain
\begin{align}
\|Z_t\|
\le
C_{\alpha_1,\alpha_2}^{(2)}
\cbr{
u^{1/\alpha_1}v_1(t)
+
u^{1/\alpha_2}v_2(t)
+
\sum_{n\ge2}R_n(t)
}.
\label{eq:mixed_finite_pointwise_common_chain}
\end{align}

Finally, by the definition of $R_n(t)$ and the change of variables
$k=n-2$, we have the exact identity
\begin{align*}
\sum_{n\ge2}R_n(t)=
\sum_{n\ge2}
\sum_{j=1}^2
2^{n/\alpha_j}d_j(t,T_{j,n-2})=
\sum_{j=1}^2
2^{2/\alpha_j}
\sum_{k\ge0}
2^{k/\alpha_j}d_j(t,T_{j,k})=
2^{2/\alpha_1}G_1(t)
+
2^{2/\alpha_2}G_2(t).
\end{align*}
Inserting this identity into
\eqref{eq:mixed_finite_pointwise_common_chain} and enlarging the
constant proves the conclusion for $u\ge4$.

We now suppose that $1\le u<4$. Apply the result already proved with
the parameter $4$. Its failure probability is $e^{-4}\le e^{-u}$.
Moreover, since $u\ge1$, by some algebra, we have
\begin{align*}
4^{1/\alpha_j}v_j(t)
\le
4^{1/\alpha_j}
u^{1/\alpha_j}v_j(t),
\ j\in\cbr{1,2}.
\end{align*}
Absorbing the constants $4^{1/\alpha_1}$ and
$4^{1/\alpha_2}$ into $A_{\alpha_1,\alpha_2}$ proves the result for
every $u\ge1$.

THus, we finish the proof.
\end{proof}

\end{document}
